\originalchapter{非线性动力学}
\originalsection{固定点与局部稳定}
自治系统$\dot x=f(x)$中, $f(x_*)=0$称固定点. 对$\eta=x-x_*$, 线性化为$\dot\eta=A\eta+O(|\eta|^2)$, $A=Df(x_*)$. 若所有特征值实部严格负, 非线性小扰动也局部渐近稳定; 可用线性稳定矩阵构造正定二次Lyapunov函数, 其导数中的负二次项支配三阶余项. 若有正实部, 存在不稳定方向. 有零或纯虚特征值时, 线性项不足以判定: 一维$\dot x=-x^3$稳定而$\dot x=x^3$不稳定, 二者在原点都有零线性化.

二维矩阵令迹为$\tau$, 行列式为$\Delta$, 特征方程$\lambda^2-\tau\lambda+\Delta=0$. $\Delta<0$是鞍点; $\Delta>0,\tau<0$为线性吸引点, $\tau>0$为排斥点; 判别式$\tau^2-4\Delta$区分节点与螺旋. $\tau=0,\Delta>0$仅给线性中心, 非线性系统是否中心仍需更高阶或守恒量判断.

\originalsection{分岔与极限环}
参数改变时固定点数目或稳定性变化称分岔. 常见一维正规形可逐一直接核对. $\dot x=\mu-x^2$在$\mu>0$有$\pm\sqrt\mu$, 导数$-2x$使正根稳定、负根不稳定, $\mu<0$无固定点, 为鞍结分岔. $\dot x=\mu x-x^2$的$x=0,\mu$交换稳定性, 为跨临界分岔. $\dot x=\mu x-x^3$在$\mu>0$出现两个稳定根$\pm\sqrt\mu$, 原点失稳, 为超临界叉形分岔. 这些分类依赖所给正规形及参数条件, 不能只看到某张图就认定所有系统都同样变化.

\begin{example}
极坐标系统$\dot r=\mu r-r^3$, $\dot\theta=\omega$, $\omega\ne0$. 求周期运动与稳定性.
\end{example}
\begin{solution}
$\mu<0$时$r>0$单调减小到0. $\mu>0$时$r=\sqrt\mu$为定半径周期轨道, 周期$2\pi/|\omega|$, 径向导数$\mu-3r^2=-2\mu<0$, 所以邻近轨道被吸引. 原点在$\mu>0$排斥. 这给出一个超临界Hopf分岔的明确模型. 自治周期轨道沿轨道的相位扰动不会衰减, 所以稳定指吸引到轨道这一集合, 不是固定相位点.
\end{solution}
平面系统满足光滑性、轨道被限制在紧区域且极限集没有固定点等条件时, Poincaré--Bendixson定理将极限行为限制为周期轨道. 本册仅准确调用这个定理, 不给其完整拓扑证明. 它不适用于三维Lorenz系统, 因而不能据平面图排除更高维混沌.

\originalsection{迭代映射与敏感依赖}
Logistic映射$x_{n+1}=rx_n(1-x_n)$在$0\le r\le4$把$[0,1]$映入自身. 固定点为$0$和$1-1/r$. 第一根导数$r$, 在$0<r<1$吸引; 第二根导数$2-r$, 在$1<r<3$吸引. $r=3$时导数到$-1$, 通常出现翻转分岔. 非固定的二周期点由$f(f(x))-x=0$去掉两个固定点因子后得
\[
x_\pm=\frac{r+1\pm\sqrt{(r-3)(r+1)}}{2r}.
\]
两个点的乘积导数为$-r^2+2r+4$, 因而该周期在$3<r<1+\sqrt6$吸引. 后续有倍周期和周期窗口, 所以不能说“所有$r>3$都混沌”.

微小扰动的线性增长因子为$\prod_{j=0}^{n-1}|f'(x_j)|$, 对数平均$\lambda=\lim_{n\to\infty}n^{-1}\sum_j\log|f'(x_j)|$称轨道Lyapunov指数（极限存在时）. 正值表示典型小误差指数放大, 但有限样本的正估计不能独自证明存在混沌吸引子. 在$r=4$时令$x=\sin^2(\pi y)$可将迭代联系到角加倍映射; 对相应不变分布的典型轨道有指数$\log2$, 特殊轨道和数值舍入仍需区别.

\originalsection{Lorenz系统与分形}
Lorenz方程为$\dot x=\sigma(y-x)$, $\dot y=x(\rho-z)-y$, $\dot z=xy-\beta z$, 其中$\sigma,\beta>0$. 散度为$-(\sigma+1+\beta)$, 因而相体积随$t$按$e^{-(\sigma+1+\beta)t}$收缩. 原点总是固定点; $\rho>1$还有$(\pm\sqrt{\beta(\rho-1)},\pm\sqrt{\beta(\rho-1)},\rho-1)$. 体积收缩意味着耗散, 并不独自证明混沌: 吸引固定点也收缩相体积. 典型参数$(10,28,8/3)$可产生复杂轨道, 本册不将数值图当全参数范围的定理.

分形可用盒维数表达: 若覆盖尺度为$\epsilon$的盒数为$N(\epsilon)$, 盒维数为$\lim_{\epsilon\to0}\log N(\epsilon)/\log(1/\epsilon)$（极限存在时）. 三等分Cantor集第$n$步有$2^n$个长$3^{-n}$的区间, 得$\log2/\log3$. 原讲稿将覆盖数公式同时称Hausdorff维数, 中文稿在此改称盒维数; 二者一般是不同定义, Cantor这个例子中数值恰好相等. 非整数维数、敏感依赖、耗散和混沌属于不同性质, 不能互相替代定义.

本章来源: 8.09第7章. 中文首版覆盖上述可逐式分析的入门结构, 不宣称完成该章的全部分岔案例、Lorenz混沌严格证明或分形测度理论.
